\documentclass[aspectratio=169]{beamer}

% ════════════════════════════════════════════════════════════════════════════
%  Tu Pana de Escritura — CTL Professional Development Conference
%  "Designing AI That Protects Student Authorship"
%
%  TIMING GUIDE
%  10-minute version : Slides 1–5, 9, 11, 13, 16, 17  (comment out 6–8, 10, 12, 14–15)
%  18-minute version : All 17 slides; expand demo pauses on slides 9 and 15
%
%  SPEAKER NOTES
%  Uncomment ONE option below, then recompile:
%    \setbeameroption{show notes on second screen=right}   <- presenter display
%    \setbeameroption{show only notes}                     <- notes-only PDF
%
%  SCREENSHOTS
%  Live screenshots loaded from:
%    /Users/Victor1/Desktop/Manual Library/Screenshots/
%  \graphicspath is set below; filenames only needed in \includegraphics calls.
% ════════════════════════════════════════════════════════════════════════════

\usepackage{pgfpages}
% \setbeameroption{show notes on second screen=right}
% \setbeameroption{show only notes}

% ── Encoding & typography ────────────────────────────────────────────────────
\usepackage[T1]{fontenc}
\usepackage[utf8]{inputenc}
\usepackage{lmodern}
\usepackage{microtype}

% ── Layout & graphics ────────────────────────────────────────────────────────
\usepackage{graphicx}
\graphicspath{{/Users/Victor1/Desktop/Manual Library/Screenshots/}}
\usepackage{booktabs}
\usepackage{colortbl}
\usepackage{array}
\usepackage{xcolor}
\usepackage{tikz}
\usepackage{setspace}
\usepackage{ragged2e}
\usepackage{hyperref}
\hypersetup{colorlinks=true, urlcolor=tpgreen, linkcolor=tpgreen}

% ── Tu Pana brand palette ────────────────────────────────────────────────────
\definecolor{tpdark}{HTML}{1a1510}
\definecolor{tpbrown}{HTML}{2b2318}
\definecolor{tpgreen}{HTML}{2d7a5f}
\definecolor{tpgold}{HTML}{c5a84f}
\definecolor{tpwarm}{HTML}{f4ede4}
\definecolor{tpred}{HTML}{c84b3b}
\definecolor{tpmuted}{HTML}{9a8e7e}
\definecolor{tpbg}{HTML}{faf8f5}
\definecolor{tpborder}{HTML}{e0d8ce}
\definecolor{tplightgreen}{HTML}{eaf4ef}

% ── Beamer theme ─────────────────────────────────────────────────────────────
\usetheme{default}
\usecolortheme{default}
\usefonttheme[onlymath]{serif}

\setbeamercolor{background canvas}{bg=tpbg}
\setbeamercolor{normal text}{fg=tpdark}
\setbeamercolor{frametitle}{fg=tpdark, bg=tpbg}
\setbeamerfont{frametitle}{size=\large, series=\bfseries}
\setbeamerfont{framesubtitle}{size=\small, series=\mdseries}
\setbeamercolor{framesubtitle}{fg=tpmuted}

\setbeamertemplate{frametitle}{%
  \vspace*{0.55em}%
  \hspace*{0.15em}%
  \begin{beamercolorbox}[wd=\textwidth, left]{frametitle}%
    {\color{tpgreen}\rule{2.6em}{2pt}}\\[0.22em]%
    \usebeamerfont{frametitle}\insertframetitle%
    \ifx\insertframesubtitle\@empty\else\\[-0.05em]%
      {\usebeamerfont{framesubtitle}\color{tpmuted}\insertframesubtitle}%
    \fi%
  \end{beamercolorbox}%
}

\setbeamercolor{itemize item}{fg=tpgreen}
\setbeamercolor{itemize subitem}{fg=tpgold}
\setbeamertemplate{itemize item}{\small\textbullet}
\setbeamertemplate{itemize subitem}{\scriptsize\textendash}
\setbeamercolor{enumerate item}{fg=tpgreen}
\setbeamerfont{enumerate item}{series=\bfseries}

\setbeamercolor{block title}{fg=tpwarm, bg=tpgreen}
\setbeamercolor{block body}{fg=tpdark, bg=tplightgreen}
\setbeamerfont{block title}{size=\small, series=\bfseries}
\setbeamerfont{block body}{size=\small}

\setbeamertemplate{navigation symbols}{}
\setbeamertemplate{footline}{%
  \vspace{0.15em}%
  \hspace{1.2em}%
  {\tiny\color{tpmuted}%
    \textit{Tu Pana de Escritura} \textperiodcentered{} CTL Professional Development%
  }%
  \hfill%
  {\tiny\color{tpmuted}\insertframenumber\,/\,\inserttotalframenumber\hspace{1.2em}}%
  \vspace{0.35em}%
}

\setbeamercolor{note page}{bg=tpwarm, fg=tpdark}
\setbeamerfont{note page}{size=\footnotesize}

% ── Macros ───────────────────────────────────────────────────────────────────

% Thin gold rule spanning text width
\newcommand{\tprule}{{\color{tpgold}\rule{\linewidth}{0.5pt}}}

% Quote callout
\newcommand{\tpspine}[1]{%
  \vspace{0.4em}
  \noindent{\color{tpgreen}\rule{1.8em}{2pt}}\\[0.25em]%
  {\itshape\color{tpdark}#1}%
  \vspace{0.4em}%
}

% ── Document metadata ────────────────────────────────────────────────────────
\title{Tu Pana de Escritura}
\subtitle{Designing AI That Protects Student Authorship}
\author{V\'{i}ctor Torres-V\'{e}lez, Ph.D.}
\institute{Department of Languages \& Linguistics \textperiodcentered{} Hostos Community College, CUNY}
\date{CTL Professional Development Conference}

% ════════════════════════════════════════════════════════════════════════════
\begin{document}
% ════════════════════════════════════════════════════════════════════════════

% ── SLIDE 1  Title ───────────────────────────────────────────────────────────
{%
\setbeamertemplate{footline}{}
\begin{frame}[plain]
  \begin{tikzpicture}[remember picture, overlay]
    \fill[tpdark]
      (current page.south west) rectangle (current page.north east);
    \fill[tpgreen]
      (current page.south west)
      rectangle ([xshift=\paperwidth, yshift=4pt] current page.south west);
    \fill[tpgold]
      (current page.south west)
      rectangle ([xshift=\paperwidth, yshift=1.8pt] current page.south west);
  \end{tikzpicture}
  \vspace{2.2em}
  \begin{center}
    {\color{tpwarm}\LARGE\bfseries Tu Pana de Escritura}\\[0.45em]
    {\color{tpgold}\large Designing AI That Protects Student Authorship}\\[1.6em]
    {\color{tpborder}\rule{14em}{0.4pt}}\\[1em]
    {\color{tpwarm}\normalsize V\'{i}ctor Torres-V\'{e}lez, Ph.D.}\\[0.25em]
    {\color{tpmuted}\small Department of Languages \& Linguistics}\\
    {\color{tpmuted}\small Hostos Community College, CUNY}\\[0.9em]
    {\color{tpmuted}\footnotesize CTL Professional Development Conference}
  \end{center}
  \note{%
    \textbf{SLIDE 1 — Opening}\\[0.4em]
    Good morning. By the end of this talk I hope you will see one possibility for using generative AI in writing instruction that does not begin with efficiency, automation, or output.\\[0.4em]
    I want to ask a different question: What would it mean to design AI around student authorship? Not around faster essays. Not around smoother prose. Not around replacing the messy, difficult, human work of writing. But around protecting that work.\\[0.4em]
    That is the question behind Tu Pana de Escritura --- a bilingual Spanish-English AI writing coach I have been developing for multilingual, Latinx, Caribbean, immigrant, and working-class college students in the Bronx.\\[0.4em]
    The central claim I want to make today is simple:\\[0.3em]
    \textbf{Tu Pana is not designed to help students avoid writing. It is designed to help students stay with writing.}\\[0.4em]
    \textit{[Pause. Let that sentence land. Then advance.]}
  }
\end{frame}
}

% ── SLIDE 2  The Problem ─────────────────────────────────────────────────────
\begin{frame}
  \frametitle{The AI Writing Problem Is Not Only Cheating}
  \framesubtitle{It is displacement}

  \medskip

  \begin{columns}[T]
    \column{0.54\textwidth}
      \setstretch{1.4}
      When students use a generic AI tool to produce a thesis, an outline, a paragraph, or an entire essay, the obvious concern is \textbf{academic integrity}.\\[0.7em]
      But the deeper problem is \textbf{displacement}.\\[0.7em]
      AI can displace the student from the very process that writing is supposed to develop:
      \begin{itemize}
        \item remembering \quad\textperiodcentered\quad questioning
        \item connecting \quad\textperiodcentered\quad structuring
        \item revising \quad\textperiodcentered\quad reflecting
      \end{itemize}
      \vspace{0.6em}
      A student may end up with a polished product\\
      but without doing the intellectual work it represents.

    \column{0.43\textwidth}
      \includegraphics[width=\linewidth, height=5.4cm, keepaspectratio]{SCR-20260519-lepw.png}
  \end{columns}

  \note{%
    \textbf{SLIDE 2 — The Problem}\\[0.4em]
    The conversation about AI and writing too often starts with cheating. That concern is real. But I want to suggest the deeper problem is displacement.\\[0.4em]
    The real question is: Can we design an AI writing environment where the student remains the author?\\[0.4em]
    \textit{[Screenshot: A student asked for help "without rewriting it" — the coach responds with a question rather than prose. The refusal is the feature. Point this out explicitly.]}%
  }
\end{frame}

% ── SLIDE 3  The Students ────────────────────────────────────────────────────
% [10-MIN: BRIEF — read the framing sentence; skip bullet list]
\begin{frame}
  \frametitle{Multilingual Students Do Not Arrive Empty-Handed}
  \framesubtitle{Tu Pana begins from a specific teaching context}

  \medskip

  \begin{columns}[T]
    \column{0.54\textwidth}
      \setstretch{1.4}
      This app emerges from teaching practice with students in the Bronx who bring:
      \begin{itemize}
        \item Multiple languages and full linguistic repertoires
        \item Migration histories and family knowledge
        \item Community memory and place-based experience
        \item Complex, generative relationships with academic institutions
      \end{itemize}
      \vspace{0.6em}
      {\color{tpred}\textbf{Generic AI can intensify the problem.}}\\[0.3em]
      It makes writing smoother --- but also more generic.\\
      More institutionally legible --- but less grounded.

    \column{0.43\textwidth}
      \includegraphics[width=\linewidth, height=5.4cm, keepaspectratio]{SCR-20260519-lcmw.png}
  \end{columns}

  \note{%
    \textbf{SLIDE 3 — The Students}\\[0.4em]
    Too often, these students receive the message that their language is a problem to be corrected, their experience merely anecdotal, their knowledge only legitimate once translated into a standardized academic voice.\\[0.4em]
    Tu Pana starts from a different premise: students already bring intellectual assets. The app's job is to help them recognize, develop, and protect those assets.\\[0.4em]
    \textit{[Screenshot: The onboarding screen displaying "Tú ya produces conocimiento" — you already produce knowledge. That sentence is the entire pedagogical claim in seven words. Let faculty read it before speaking.]}%
  }
\end{frame}

% ── SLIDE 4  What Tu Pana Is ─────────────────────────────────────────────────
\begin{frame}
  \frametitle{A Bilingual AI Writing Coach}
  \framesubtitle{Not an essay generator \textperiodcentered{} not a chatbot \textperiodcentered{} not a productivity tool}

  \medskip

  \begin{columns}[T]
    \column{0.50\textwidth}
      \textit{Tu Pana de Escritura} means \textbf{``Your Writing Companion.''}\\[0.7em]
      The most important thing to understand is not what the AI does.\\[0.4em]
      \textbf{It is what the AI is not allowed to do.}

      \vspace{0.7em}
      {\color{tpred}The AI does not:}
      \begin{itemize}
        \item write the essay
        \item generate the thesis or outline
        \item produce the introduction or conclusion
        \item invent citations
        \item write the student's self-assessment
      \end{itemize}

      \vspace{0.5em}
      The AI asks. It prompts. It supports.\\
      \textbf{The writing remains the student's.}

    \column{0.47\textwidth}
      \includegraphics[width=\linewidth, height=5.6cm, keepaspectratio]{SCR-20260519-ledl.png}

      \vspace{0.5em}
      \begin{block}{Scaffolded Authorship}
        \centering\small
        Supporting the student's development as an author\\
        through structured, constrained AI interaction.
      \end{block}
  \end{columns}

  \note{%
    \textbf{SLIDE 4 — What Tu Pana Is}\\[0.4em]
    Tu Pana de Escritura is a bilingual Spanish-English web app that guides students through a ten-stage writing process.\\[0.4em]
    The most important thing to understand is not what the AI does. It is what the AI is not allowed to do.\\[0.4em]
    This is what I call scaffolded authorship.\\[0.4em]
    \textit{[Screenshot: Main interface with all ten stage names visible across the navigation bar. The stages themselves make the argument — this is a sequence, not a chatbot. Let the audience read the stage names.]}%
  }
\end{frame}

% ── SLIDE 5  Student-Writes-First ────────────────────────────────────────────
\begin{frame}
  \frametitle{The Student Writes First}
  \framesubtitle{The AI responds second --- always}

  \medskip

  \begin{columns}[T]
    \column{0.54\textwidth}
      \setstretch{1.4}
      At every meaningful stage, the student must write \textbf{before} the AI can respond.\\[0.7em]
      This matters because writing is not just the production of text.\\[0.4em]
      \textbf{Writing is a way of thinking.}\\[0.7em]
      When students write, they make choices:
      \begin{itemize}\footnotesize
        \item What do I remember? \; What matters?
        \item What evidence do I trust?
        \item How do I want to sound?
        \item What do I want the reader to understand?
      \end{itemize}
      \vspace{0.5em}
      If the AI answers those questions first, the student receives text --- but loses the practice of \textit{thinking through writing}.

    \column{0.43\textwidth}
      \includegraphics[width=\linewidth, height=5.6cm, keepaspectratio]{SCR-20260519-lemu.png}

      \vspace{0.5em}
      \tprule
      \vspace{0.35em}
      \begin{center}
        {\small\itshape The student writes. The AI asks.\\The student decides.}\\[0.2em]
        {\footnotesize\color{tpmuted}That is the ethical architecture of the tool.}
      \end{center}
  \end{columns}

  \note{%
    \textbf{SLIDE 5 — Student-Writes-First Principle}\\[0.4em]
    Tu Pana places AI in a secondary role. The student writes. The AI asks. The student decides. That is the ethical architecture of the tool.\\[0.4em]
    Return to the spine of the talk here: Tu Pana is not designed to help students avoid writing. It is designed to help students stay with writing.\\[0.4em]
    \textit{[Screenshot: Student draft in the text area with the "Enviar" send button visible. The input box is the dominant UI element — the student's text is what the app is waiting for. That layout is the student-writes-first principle made visible.]}%
  }
\end{frame}

% ── SLIDE 6  Ten Stages ──────────────────────────────────────────────────────
% [10-MIN: SKIP — add one sentence to Slide 4: "The app has ten stages,
%  from personal memory to process reflection, and moves through them in sequence."]
\begin{frame}
  \frametitle{Ten Stages: From Memory to Reflection}
  \framesubtitle{A structured sequence, not an open-ended prompt}

  \medskip

  \begin{columns}[T]
    \column{0.47\textwidth}
      \footnotesize\setstretch{1.35}
      \begin{enumerate}
        \item \textbf{An\'{e}cdota / Anecdote}\\{\color{tpmuted}personal or family memory}
        \item \textbf{Conexi\'{o}n / Connection}\\{\color{tpmuted}social \& historical context}
        \item \textbf{Tu Pitch / Topic Pitch}\\{\color{tpmuted}working claim}
        \item \textbf{Investigaci\'{o}n / Research}\\{\color{tpmuted}sources, evidence, credibility}
        \item \textbf{Esquema / Outline}\\{\color{tpmuted}student-generated structure}
      \end{enumerate}

    \column{0.47\textwidth}
      \footnotesize\setstretch{1.35}
      \begin{enumerate}\setcounter{enumi}{5}
        \item \textbf{Primer Borrador / First Draft}\\{\color{tpmuted}full student draft}
        \item \textbf{Revisi\'{o}n / Revision}\\{\color{tpmuted}Five Questions framework}
        \item \textbf{Pulir Voz / Voice Polish}\\{\color{tpmuted}Voice Vault}
        \item \textbf{Checklist}\\{\color{tpmuted}metacognitive self-evaluation}
        \item \textbf{Mi Cierre de Proceso}\\{\color{tpmuted}writing snapshot \& report}
      \end{enumerate}
  \end{columns}

  \vspace{0.9em}
  \begin{center}
    {\small\color{tpgreen}\itshape
      memory $\;\to\;$ connection $\;\to\;$ research $\;\to\;$ draft
      $\;\to\;$ revision $\;\to\;$ voice $\;\to\;$ reflection}
  \end{center}

  \note{%
    \textbf{SLIDE 6 — Ten-Stage Process}\\[0.4em]
    The movement is intentional: memory to reflection. The app is not helping students finish an assignment. It is teaching them a process.\\[0.4em]
    Worth pausing on the opening: it is unusual to begin academic writing with anecdote --- to ask students to write from what they know before asking them to research. That inversion is deliberate. It establishes the student as a narrator and a witness before asking them to be a scholar.\\[0.4em]
    \textit{[18-minute version: Walk through all ten stages with one sentence each. Pause on Stage 1 (memory as starting point) and Stage 10 (reflection as closing). Those two stages anchor the arc.]}\\[0.3em]
    \textit{[10-minute version: Skip this slide entirely.]}
  }
\end{frame}

% ── SLIDE 7  Decolonial AI Pedagogy ─────────────────────────────────────────
% [10-MIN: SKIP — name the concept once in Slide 3 and move on]
\begin{frame}
  \frametitle{Decolonial AI Pedagogy}
  \framesubtitle{Redesigning the role of AI --- not just the content of AI}

  \medskip

  \begin{columns}[T]
    \column{0.54\textwidth}
      \setstretch{1.4}
      Decolonial AI does \textbf{not} mean adding diverse content to a chatbot.\\[0.5em]
      It means \textbf{redesigning the role of AI} in the educational encounter.\\[0.7em]
      The deeper questions:
      \begin{itemize}
        \item Whose thinking does the AI scaffold?
        \item Whose labor does it displace?
        \item Whose language does it normalize?
        \item Whose knowledge does it treat as valid?
        \item Whose voice does it flatten?
      \end{itemize}
      \vspace{0.5em}
      Tu Pana constrains AI so it cannot become another mechanism for
      \textbf{extraction, standardization, or ghostwriting}.

    \column{0.43\textwidth}
      \includegraphics[width=\linewidth, height=5.6cm, keepaspectratio]{SCR-20260519-libi.png}

      \vspace{0.5em}
      \begin{block}{}
        \centering\small
        Not the absence of AI.\\
        The \textbf{placement} of AI in the service\\
        of student intellectual sovereignty.
      \end{block}
  \end{columns}

  \note{%
    \textbf{SLIDE 7 — Decolonial AI Pedagogy}\\[0.4em]
    A decolonial question is not only: Does the tool include Spanish? Does it mention Latinx students? Does it sound inclusive?\\[0.4em]
    The deeper questions are about power, labor, language, and voice --- and Tu Pana tries to answer those questions through design.\\[0.4em]
    \textit{[Screenshot: Stage 4 AI literacy checkpoint — the app asking students to evaluate how AI handles information and uncertainty. The tool turns itself into an object of critical inquiry. That is what decolonial AI pedagogy looks like in practice.]}%
  }
\end{frame}

% ── SLIDE 8  Bilingualism \& Voice ───────────────────────────────────────────
% [10-MIN: SKIP — covered by slide 9's Voice Vault framing]
\begin{frame}
  \frametitle{Bilingualism, Translanguaging, and Voice}
  \framesubtitle{Spanglish is not a glitch. It can be a rhetorical choice.}

  \medskip

  \begin{columns}[T]
    \column{0.54\textwidth}
      \setstretch{1.4}
      Tu Pana treats Spanish, English, Spanglish, code-switching, dialect, accent, family phrases, and community idioms as \textbf{rhetorical resources} --- not as errors.\\[0.7em]
      When a student writes \textit{``mi abuela me dec\'{i}a''} in an English-language essay, the app does not correct it.\\[0.5em]
      Instead it asks:\\[0.3em]
      {\footnotesize\itshape Why does this phrase matter here?\\
       What would be lost if you translated it?\\
       Is this a phrase you want to protect?}\\[0.7em]
       This is \textbf{multilingual rhetorical agency}: not the suppression of the student's full linguistic repertoire, but the development of their capacity to make \textit{intentional choices} about how to deploy it.

    \column{0.43\textwidth}
      \includegraphics[width=\linewidth, height=5.8cm, keepaspectratio]{SCR-20260519-lfbj.png}
  \end{columns}

  \note{%
    \textbf{SLIDE 8 — Bilingualism and Voice}\\[0.4em]
    The goal is not to tell students: write however you want and audience does not matter. The goal is to help students develop judgment about their own language choices.\\[0.4em]
    This is what multilingual rhetorical agency looks like in practice.\\[0.4em]
    \textit{[Screenshot: Coach response in Spanish — fully bilingual, not translated. The AI's language choice signals whose students this is for. The behavior of the AI is the argument.]}%
  }
\end{frame}

% ── SLIDE 9  Voice Vault ─────────────────────────────────────────────────────
\begin{frame}
  \frametitle{The Voice Vault}
  \framesubtitle{Revision without erasure \textperiodcentered{} Stage~8: Pulir Voz / Voice Polish}

  \medskip

  \begin{columns}[T]
    \column{0.49\textwidth}
      \setstretch{1.4}
      Students identify phrases in their draft that must be \textbf{protected during revision}:
      \begin{itemize}
        \item A word from a parent or grandparent
        \item A code-switched phrase
        \item A community-specific term
        \item A sentence that sounds \textit{exactly like them}
        \item A moment where translation would weaken the meaning
      \end{itemize}
      \vspace{0.6em}
      The student declares:\\[0.3em]
      {\itshape\small ``This language is deliberate.\\
       This is not an error.\\
       Do not flatten this.''}

      \vspace{0.7em}
      \tprule
      \vspace{0.3em}
      {\small\color{tpgreen}\bfseries Smoother is not always better.}

    \column{0.48\textwidth}
      \includegraphics[width=\linewidth, height=6.8cm, keepaspectratio]{SCR-20260519-lmic.png}

      \vspace{0.4em}
      {\footnotesize\color{tpmuted}\itshape
       A smoother sentence can be less precise.\\
       A more polished sentence can be less honest.\\
       A more standard sentence can carry less cultural knowledge.}
  \end{columns}

  \note{%
    \textbf{SLIDE 9 — Voice Vault}\\[0.4em]
    Stage 8 includes one of the features I care about most. The Voice Vault exists because AI revision tools often push writing toward a generic version of academic clarity. Tu Pana teaches students that revision should strengthen their writing without erasing them from it.\\[0.4em]
    Pause on the phrase: smoother is not always better. Let it sit. Faculty will recognize its implications for how they think about academic writing standards.\\[0.4em]
    \textit{[Screenshot: Stage 8 Voice Vault — the interface showing the transition into voice-protection mode. Let faculty look at it before you speak. The design communicates that revision is not standardization.]}\\[0.3em]
    \textit{[10-minute version: This slide carries the entire bilingualism argument. Say "smoother is not always better" and move on.]}%
  }
\end{frame}

% ── SLIDE 10  Research \& Citation Ethics ────────────────────────────────────
% [10-MIN: SKIP]
\begin{frame}
  \frametitle{Research and Citation Ethics}
  \framesubtitle{The AI does not invent sources --- research is part of the learning}

  \medskip

  \begin{columns}[T]
    \column{0.54\textwidth}
      \setstretch{1.4}
      One of the known risks of generative AI: it can produce \textbf{fabricated citations} that sound plausible but do not exist.\\[0.7em]
      In Tu Pana, when students need sources, the AI \textit{redirects} them toward the research process:
      \begin{itemize}
        \item What question are you trying to answer?
        \item What keywords might help?
        \item What kinds of sources count as evidence?
        \item How does this source connect to your argument?
        \item How will you know if a source is credible?
      \end{itemize}

    \column{0.43\textwidth}
      \includegraphics[width=\linewidth, height=5.6cm, keepaspectratio]{SCR-20260519-lilq.png}

      \vspace{0.5em}
      \begin{block}{}
        \centering\small
        Students must locate, read, evaluate,\\
        and cite \textbf{real sources} themselves.\\[0.2em]
        {\footnotesize This is not a limitation.\\It is a pedagogical commitment.}
      \end{block}
  \end{columns}

  \note{%
    \textbf{SLIDE 10 — Research and Citation Ethics}\\[0.4em]
    Tu Pana does not generate citations. In a writing class, fabricated sources are not a minor technical issue --- they undermine the integrity of research as an intellectual practice.\\[0.4em]
    The AI can support research thinking. But students must locate, read, evaluate, and cite real sources. Again: the app protects the process.\\[0.4em]
    \textit{[Screenshot: Research stage guardrail notice — the app explicitly declining to generate citations and redirecting the student toward the research process. Read the guardrail text aloud if time permits.]}%
  }
\end{frame}

% ── SLIDE 11  Five Questions Framework ───────────────────────────────────────
\begin{frame}
  \frametitle{The Five Questions Framework}
  \framesubtitle{A practical framework for student self-reading}

  \medskip

  \begin{columns}[T]
    \column{0.49\textwidth}
      \setstretch{1.45}
      Students evaluate their writing --- and \textbf{the AI's feedback} --- using five questions:

      \vspace{0.5em}
      \begin{enumerate}
        \item Is it \textbf{accurate}?
        \item Does it \textbf{preserve voice}?
        \item Is it \textbf{specific}?
        \item Does it \textbf{show thinking}?
        \item Does it \textbf{honor cultural knowledge}?
      \end{enumerate}

      \vspace{0.7em}
      {\small\color{tpmuted}These questions also apply to the AI's feedback.}\\[0.2em]
      \textbf{The AI is not the authority.}\\[0.2em]
      {\small Its feedback is something to question, accept, revise, or reject.\\
       That is part of AI literacy.}

    \column{0.48\textwidth}
      \includegraphics[width=\linewidth, height=5.8cm, keepaspectratio]{SCR-20260519-ldzz.png}

      \vspace{0.4em}
      {\footnotesize\color{tpgreen}\itshape
       Not grammar. Not correctness.\\
       Groundedness, voice, specificity,\\
       thought, and cultural accountability.}
  \end{columns}

  \note{%
    \textbf{SLIDE 11 — Five Questions Framework}\\[0.4em]
    These are not rubric categories. They are habits of critical self-reading. They teach students that good writing is not just writing that sounds polished.\\[0.4em]
    Pause on question 5: Does it honor cultural knowledge? That one lands hardest with faculty who teach multilingual students. It reframes what we mean by strong writing.\\[0.4em]
    \textit{[Screenshot: All Five Questions displayed together as a summary panel. Read Question 5 — "Does it honor cultural knowledge?" — aloud and let it sit. Faculty already value this kind of critical thinking; it just hasn't been applied to writing and AI before.]}\\[0.3em]
    \textit{[10-minute version: Read all five questions aloud. Pause on \#5. Move on.]}%
  }
\end{frame}

% ── SLIDE 12  Process Documentation ─────────────────────────────────────────
% [10-MIN: One sentence only: "The app ends with an Instructor Process Report ---
%  not a grade, just a record of the student's process."]
\begin{frame}
  \frametitle{In the Age of AI, Process Becomes Evidence}
  \framesubtitle{Mi Cierre de Proceso \textperiodcentered{} My Writing Snapshot \textperiodcentered{} Instructor Process Report}

  \medskip

  \begin{columns}[T]
    \column{0.52\textwidth}
      \setstretch{1.4}
      The final essay no longer tells us enough.\\[0.7em]
      At Stage~10, students complete \textbf{Mi Cierre de Proceso}:
      \begin{itemize}
        \item What did I learn?
        \item What was most difficult?
        \item How did my thinking change?
        \item What would I do differently?
      \end{itemize}
      \vspace{0.5em}
      \textit{Only after the student writes} can the AI offer a limited coach perspective --- based on what it observed in the app.\\[0.5em]
      The \textbf{Instructor Process Report} is then generated.

      \vspace{0.4em}
      \tprule
      \vspace{0.3em}
      {\small Not a grade. Not an evaluation.\\
       \textbf{Process documentation.}}

    \column{0.45\textwidth}
      \includegraphics[width=\linewidth, height=6.5cm, keepaspectratio]{SCR-20260519-lnae.png}
  \end{columns}

  \note{%
    \textbf{SLIDE 12 — Process Documentation}\\[0.4em]
    Tu Pana responds to the academic integrity challenge not by trying to detect AI use after the fact, but by making the process visible from the beginning.\\[0.4em]
    The Instructor Process Report helps instructors see the student's journey, not just the final product. It is meant to supplement instructor judgment, not replace it.\\[0.4em]
    \textit{[Screenshot: Stage 10 Mi Cierre de Proceso completion modal — "Llegaste al final con dignidad · You reached the end with dignity." Below it: the self-assessment form with categories including Voice/Voz, AI Judgment/Criterio sobre la IA, Cultural Knowledge. That last category does not exist in any standard rubric. Note it explicitly.]}%
  }
\end{frame}

% ── SLIDE 13  Why This Matters for Faculty ───────────────────────────────────
\begin{frame}
  \frametitle{Why This Matters for Faculty}
  \framesubtitle{A response to pressures faculty already feel}

  \medskip

  \begin{columns}[T]
    \column{0.52\textwidth}
      \setstretch{1.4}
      \begin{itemize}
        \item \textbf{Post-pandemic writing challenges}\\
          {\small\color{tpmuted}Students need more scaffolding and process support}
        \item \textbf{Multilingual writing anxiety}\\
          {\small\color{tpmuted}Students arrive fearing their language is wrong before they begin}
        \item \textbf{AI ghostwriting}\\
          {\small\color{tpmuted}Faculty need process evidence, not just detection tools}
        \item \textbf{AI literacy}\\
          {\small\color{tpmuted}Students need to learn to question AI, not just use it}
        \item \textbf{Equity}\\
          {\small\color{tpmuted}AI tools need to be built for students who need the most support}
      \end{itemize}

    \column{0.45\textwidth}
      \includegraphics[width=\linewidth, height=5.0cm, keepaspectratio]{SCR-20260519-lkxv.png}

      \vspace{0.5em}
      \begin{block}{}
        \centering\small
        Tu Pana does not pretend AI doesn't exist.\\
        It does not rely on surveillance.\\
        It does not outsource faculty judgment.
      \end{block}
  \end{columns}

  \note{%
    \textbf{SLIDE 13 — Why This Matters for Faculty}\\[0.4em]
    For faculty, Tu Pana responds to several pressures at once. Read through the five bullets, pausing briefly on each.\\[0.4em]
    Instead of panic, prohibition, or detection, Tu Pana asks: Can AI be placed inside a pedagogical structure that protects student authorship? That is the experiment.\\[0.4em]
    \textit{[Screenshot: Stage 7 view showing the complete stage navigation bar across the top of the app — all stages visible, student's current progress visible. The full arc of the process is readable at a glance. This answers the question: "Is this just ChatGPT?" visually and immediately.]}%
  }
\end{frame}

% ── SLIDE 14  Comparison Table ───────────────────────────────────────────────
% [10-MIN: SKIP — cover in one sentence during Slide 4]
\begin{frame}
  \frametitle{Generic AI Produces Outputs}
  \framesubtitle{Tu Pana develops authorship}

  \medskip

  {\footnotesize
    \setlength{\tabcolsep}{0.55em}
    \renewcommand{\arraystretch}{1.28}
    \begin{tabular}{>{\bfseries}p{3.0cm} p{4.2cm} p{4.6cm}}
      \toprule
      \rowcolor{tplightgreen}
      & {\normalfont\bfseries\color{tpmuted}Generic AI (no guardrails)}
      & {\normalfont\bfseries\color{tpgreen}Tu Pana de Escritura} \\
      \midrule
      Process structure
        & None --- output on demand
        & 10 sequential scaffolded stages \\
      Authorship
        & Can generate full essays \& theses
        & Constrained from generating essay prose \\
      Language policy
        & Defaults to standardized English
        & Treats bilingualism as a resource \\
      Citation integrity
        & Can fabricate sources
        & Will not generate citations \\
      Revision approach
        & May flatten voice
        & Revision without erasure; Voice Vault \\
      Process evidence
        & None
        & Instructor Process Report \\
      AI literacy
        & Not addressed
        & Embedded checkpoints at 3 stages \\
      Designed for
        & General market
        & Multilingual, Latinx, immigrant students \\
      \bottomrule
    \end{tabular}
  }

  \note{%
    \textbf{SLIDE 14 — Generic AI vs.\ Tu Pana}\\[0.4em]
    Tu Pana is not ChatGPT with a friendly Spanish name. The difference is not branding. The difference is architecture.\\[0.4em]
    Walk through the table. Do not rush. Faculty skeptics will find the rows on authorship, language policy, citation integrity, and AI literacy most persuasive.\\[0.4em]
    End by returning to the spine: Generic AI is product-centered. Tu Pana is process-centered.\\[0.4em]
    \textbf{Tu Pana is not designed to help students avoid writing. It is designed to help students stay with writing.}
  }
\end{frame}

% ── SLIDE 15  Demo Moment ────────────────────────────────────────────────────
\begin{frame}
  \frametitle{What the AI Says --- and What It Refuses to Say --- Matters}
  \framesubtitle{A brief demonstration}

  \medskip

  \begin{columns}[T, totalwidth=\textwidth]
    \column{0.49\textwidth}
      \includegraphics[width=\linewidth, height=3.7cm, keepaspectratio]{SCR-20260519-lemu.png}
      \vspace{0.35em}
      \includegraphics[width=\linewidth, height=3.7cm, keepaspectratio]{SCR-20260519-lfbj.png}

    \column{0.49\textwidth}
      \includegraphics[width=\linewidth, height=3.7cm, keepaspectratio]{SCR-20260519-lepw.png}
      \vspace{0.35em}
      \includegraphics[width=\linewidth, height=3.7cm, keepaspectratio]{SCR-20260519-lmox.png}
  \end{columns}

  \vspace{0.5em}
  \begin{center}
    {\small\color{tpgreen}\itshape
      In educational AI, refusal is not always failure.\\
      Sometimes refusal is pedagogy.}
  \end{center}

  \note{%
    \textbf{SLIDE 15 — Demo Moment}\\[0.4em]
    What I want the audience to notice is not only the interface. I want them to notice the behavior.\\[0.4em]
    When a student asks for help, the coach does not take over. It asks questions. It gives structure without filling it in. It redirects requests for ghostwriting. It protects research integrity. It treats bilingual language as meaningful.\\[0.4em]
    The most important feature of Tu Pana may be what it refuses to do.\\[0.4em]
    \textit{[2x2 grid: Top-left — student draft in input area (student writes first). Top-right — coach response in Spanish, asking questions not generating prose. Bottom-left — student asking for help "without rewriting it"; coach refuses and asks instead. Bottom-right — stage progression view, student advancing through the process.]}\\[0.3em]
    \textit{[18-minute version: Pause on each screenshot. Ask: "What do you notice?" Give 30 seconds before explaining.]}%
  }
\end{frame}

% ── SLIDE 16  The Larger Claim ───────────────────────────────────────────────
\begin{frame}
  \frametitle{We Are Not Just Choosing Tools}
  \framesubtitle{We are choosing pedagogical values}

  \medskip

  \begin{columns}[T]
    \column{0.54\textwidth}
      \setstretch{1.4}
      The field of AI-assisted writing instruction is still being built.\\[0.7em]
      \textbf{The choices we make now matter.}\\[0.7em]
      Will AI become a shortcut around writing --- a machine for producing acceptable academic prose?\\[0.7em]
      \textbf{Or can AI be designed differently?}
      \begin{itemize}
        \item Can it slow students down productively?
        \item Can it protect multilingual voice?
        \item Can it help faculty see \textit{process} rather than only \textit{product}?
        \item Can it be built around equity, not automation?
      \end{itemize}

    \column{0.43\textwidth}
      \includegraphics[width=\linewidth, height=5.6cm, keepaspectratio]{SCR-20260519-lnae.png}

      \vspace{0.5em}
      \begin{block}{}
        \centering\small
        Tu Pana is one attempt to answer \textbf{yes}.\\[0.2em]
        {\footnotesize\color{tpdark!70}Not perfectly. Not finally. But deliberately.}
      \end{block}
  \end{columns}

  \note{%
    \textbf{SLIDE 16 — The Larger Claim}\\[0.4em]
    We are deciding what AI is allowed to do in the writing classroom. Those decisions have consequences --- for whose thinking gets scaffolded, whose voice gets protected, and whose labor gets protected or displaced.\\[0.4em]
    Tu Pana is one answer to those questions. It is offered not as the only answer, but as proof that the questions can be asked with care, and that the answers can be built with integrity.\\[0.4em]
    \textit{[Screenshot: Stage 10 completion modal — "Llegaste al final con dignidad." End with dignity. That phrase is the answer to every question on this slide. This is what deliberate design looks like.]}%
  }
\end{frame}

% ── SLIDE 17  Contributions ──────────────────────────────────────────────────
{%
\setbeamertemplate{footline}{}
\begin{frame}
  \frametitle{Contributions}
  \framesubtitle{What Tu Pana de Escritura offers}

  \medskip

  \begin{columns}[T]
    \column{0.57\textwidth}
      \setstretch{1.5}
      \begin{enumerate}
        \item A model of \textbf{guardrailed AI} for writing instruction that prevents essay generation while supporting students
        \item A model of \textbf{decolonial AI pedagogy} designed around student dignity, authorship, and linguistic justice
        \item A model of \textbf{multilingual writing support} that treats Spanish, English, Spanglish, and community language as rhetorical resources
        \item A model of \textbf{process-based authorship documentation} that gives faculty a richer view of student writing development
        \item A practical \textbf{classroom response to AI} that moves beyond panic, prohibition, or detection
      \end{enumerate}

    \column{0.40\textwidth}
      \includegraphics[width=\linewidth, height=7.0cm, keepaspectratio]{SCR-20260519-lcmw.png}

      \vspace{0.5em}
      \begin{center}
        {\footnotesize\color{tpgreen}
          \href{https://vikthor1.github.io/tu-pana/}{vikthor1.github.io/tu-pana}}
      \end{center}
  \end{columns}

  \note{%
    \textbf{SLIDE 17 — Contributions (Closing)}\\[0.4em]
    Read through the five contributions. Pause on each one.\\[0.4em]
    Then close with this sentence --- speak slowly, do not rush it:\\[0.4em]
    \textbf{``In a moment when AI can generate text instantly, Tu Pana insists that the student's struggle to write, remember, connect, revise, and reflect is not an obstacle to learning. It is the learning.''}\\[0.4em]
    \textit{[Screenshot behind the five contributions: the onboarding screen — "Tú ya produces conocimiento." The closing image mirrors the opening premise: the student already produces knowledge. The talk ends where Tu Pana begins.]}\\[0.4em]
    \textit{[Pause. Do not say Thank You. Do not say Questions. Let the closing sentence stand. Then wait for the room.]}%
  }
\end{frame}
}

% ════════════════════════════════════════════════════════════════════════════
%  TIMING REFERENCE  (not a slide)
%
%  10-MINUTE VERSION  — ~10 slides
%  ─────────────────────────────────────────────────────────────────────────
%  Keep  : 1, 2, 3 (brief), 4, 5, 9, 11, 13, 16, 17
%  Skip  : 6, 7, 8, 10, 12, 14, 15
%  Note  : On slide 4 add: "The app has ten stages — from personal memory
%           to process reflection — and moves through them in sequence."
%           On slide 12, replace with one sentence spoken aloud.
%  Pace  : ~60 seconds per slide on average.
%
%  18-MINUTE VERSION  — all 17 slides
%  ─────────────────────────────────────────────────────────────────────────
%  Expand : Slide 9 (Voice Vault) — 90 sec; pause before explaining.
%           Slide 15 (demo) — 90–120 sec; ask audience "what do you notice?"
%           Slide 6 (ten stages) — walk through all 10 with one sentence each.
%           Slide 14 (table) — let faculty read; invite one question.
%  Reserve: Final 2 minutes for slide 17 and the closing spoken line.
% ════════════════════════════════════════════════════════════════════════════

\end{document}
